\subsection{Testbed Description}
\label{subsec:testbed}

All experiments reported in this paper are conducted on a dedicated
Kubernetes cluster deployed on the OneLab experimental facility,
ensuring full reproducibility and isolation from production traffic.
The testbed is designed to emulate a federated edge-cloud deployment
at a scale representative of a single administrative domain, while
providing the instrumentation granularity required for nanosecond-level
performance characterisation.

\subsubsection{Hardware Configuration}
\label{subsubsec:hardware}

The cluster comprises six virtualised nodes hosted on the OneLab
bare-metal infrastructure, each provisioned as a KVM-based virtual
machine with dedicated CPU pinning and memory reservation to minimise
hypervisor-induced jitter.
Table~\ref{tab:hardware} summarises the hardware specification of
each node.

\begin{table}[htbp]
  \centering
  \caption{Testbed node hardware specifications.}
  \label{tab:hardware}
  \resizebox{\columnwidth}{!}{%
  \begin{tabular}{lccccc}
    \toprule
    \textbf{Node} & \textbf{Role} & \textbf{vCPUs} & \textbf{RAM (GiB)} & \textbf{NICs} & \textbf{Kernel} \\
    \midrule
    \texttt{nemo-dev-master}  & Control plane & 8  & 14.5 & 2$\times$virtio & 6.8.0-100 \\
    \texttt{nemo-dev-worker1} & Worker        & 12 & 29.2 & 2$\times$virtio & 6.8.0-100 \\
    \texttt{nemo-dev-worker2} & Worker        & 12 & 29.2 & 2$\times$virtio & 6.8.0-100 \\
    \texttt{nemo-dev-worker3} & Worker        & 12 & 29.2 & 2$\times$virtio & 6.8.0-100 \\
    \texttt{nemo-dev-worker4} & Worker        & 12 & 29.2 & 2$\times$virtio & 6.8.0-100 \\
    \texttt{nemo-dev-worker5} & Worker        & 10 & 29.2 & 2$\times$virtio & 6.8.0-79  \\
    \bottomrule
  \end{tabular}%
  }
\end{table}

All nodes are equipped with Intel Xeon Cascadelake processors
operating at a nominal frequency of 2.89~GHz, with 16~MB of shared
L3 cache per socket.
Each node exposes two paravirtualised network interfaces
(\texttt{virtio\_net} driver):
\texttt{enp1s0}, connected to the public-facing network
(\texttt{132.227.122.0/24}), and
\texttt{enp7s0}, connected to the cluster-internal management and
data plane network (\texttt{192.168.111.0/24}).
Both interfaces operate at a virtual line rate equivalent to 10~Gbps
with an MTU of 1500~bytes (1450~bytes for overlay traffic due to
VXLAN encapsulation overhead).
The operating system on all nodes is Ubuntu~24.04.2~LTS with
Linux kernel~6.8.0 (generic variant).

\subsubsection{Network Topology}
\label{subsubsec:topology}

\paragraph{Layer~2/Layer~3 underlay.}
The cluster-internal network uses a flat Layer~2 segment on the
\texttt{192.168.111.0/24} subnet, interconnecting all six nodes
via a shared virtual bridge on the OneLab infrastructure.
The public network (\texttt{132.227.122.0/24}) provides external
API access to the Kubernetes control plane at
\texttt{api.main.nemo.onelab.eu:6443} and hosts the MetalLB load
balancer pool (\texttt{132.227.122.26--30}) for externally
reachable services.
Default routing on all nodes egresses via \texttt{132.227.122.1}
on the public interface.

\paragraph{Overlay network (CNI).}
Pod-to-pod communication is managed by Flannel~v0.26.4 operating
in VXLAN mode with a pod CIDR of \texttt{10.244.0.0/16}, allocating
a \texttt{/24} subnet per node.
The VXLAN tunnel interface (\texttt{flannel.1}) encapsulates pod
traffic over the cluster-internal \texttt{192.168.111.0/24}
network, adding a 50-byte overhead (outer UDP + VXLAN header) to
each frame and resulting in an effective overlay MTU of 1450~bytes.
Multus~CNI (snapshot release) is deployed as a secondary CNI plugin
to enable attachment of additional network interfaces to pods, as
required by the L2S-M overlay segments provisioned by the mNCC.

\paragraph{Layer~2 Slice Manager (L2S-M) overlay.}
The mNCC provisions on-demand Layer~2 overlay segments via the
L2S-M Connectivity Controller (image: \texttt{alexdecb/l2sm-md:0.4}),
which exposes a gRPC interface on port~50051
(\texttt{l2sm-grpc-service.nemo-net.svc.cluster.local}).
Overlay segments are instantiated as VXLAN tunnels between
designated Network Edge Devices (NEDs) running on the worker nodes,
with segment membership managed via Network Attachment Definitions
created by the L2S-M operator.
A dedicated CoreDNS instance
(\texttt{l2sm-coredns}) resolves overlay-internal names using the
\texttt{.inter.l2sm} suffix.

\paragraph{Inter-cluster connectivity.}
Cross-cluster overlay reachability is achieved through BGP-based
underlay routing (assumption~A1).
Each cluster boundary is managed by an ExaBGP instance advertising
the local pod CIDR prefixes to neighbouring Autonomous Systems.
The ONOS SDN controller orchestrates flow rules on the NED
switches to steer cross-cluster traffic through the appropriate
VXLAN tunnels.

\subsubsection{Software Stack}
\label{subsubsec:software_stack}

Table~\ref{tab:software} details the key software components and
their versions as deployed on the testbed.

\begin{table}[htbp]
  \centering
  \caption{Software stack deployed on the evaluation testbed.}
  \label{tab:software}
  \begin{tabular}{ll}
    \toprule
    \textbf{Component} & \textbf{Version} \\
    \midrule
    Kubernetes                 & v1.29.13 \\
    Container runtime          & containerd~1.7.28 \\
    Flannel (CNI)              & v0.26.4 (VXLAN backend) \\
    Multus CNI                 & snapshot \\
    mNCC Intent-Based System   & v0.0.19 \\
    L2S-M Connectivity Controller & 0.4 \\
    NeMeX (Network Metrics)    & latest \\
    Network Probe              & v1.0.0 \\
    Prometheus Operator        & (kube-prometheus-stack) \\
    Grafana                    & (bundled) \\
    Kepler (energy telemetry)  & 0.7.10 \\
    MetalLB                    & (L2 mode) \\
    Kong API Gateway           & (nemo-sec) \\
    RabbitMQ                   & (single-node, nemo-sec) \\
    Rook-Ceph (storage)        & (RBD + CephFS) \\
    \bottomrule
  \end{tabular}
\end{table}

\subsubsection{Time Synchronisation}
\label{subsubsec:ptp}

Accurate cross-node time synchronisation is a prerequisite for
meaningful per-stage provisioning decomposition and fault detection
latency measurements.
The virtualised environment does not expose hardware PTP
(IEEE~1588) clock devices at the guest level; therefore, time
synchronisation is enforced via the \texttt{kvm-clock} clocksource
provided by the hypervisor, which synchronises all guest clocks to
the host's hardware reference.
Within each node, high-resolution timestamps are obtained using
\texttt{clock\_gettime(CLOCK\_MONOTONIC)} with nanosecond
precision, ensuring monotonicity and immunity to NTP step
adjustments.

Prior to each measurement session, inter-node clock offset is
verified by exchanging ICMP timestamp probes between all node pairs
and computing the one-way delay asymmetry.
The measured inter-node offset was consistently below 10~\textmu s
across all pairs during the evaluation period---well within the
tolerance required for the millisecond-scale measurements reported
in this work.
For the self-healing experiments (Section~\ref{subsec:selfhealing_methodology}),
where detection latencies can reach the sub-second range, this
synchronisation precision introduces a maximum measurement
uncertainty of $<\!0.01\%$ relative to the observed values.

\subsubsection{Instrumentation and Monitoring}
\label{subsubsec:instrumentation}

The testbed deploys a comprehensive monitoring stack to provide
continuous observability of network, compute, and energy metrics:

\begin{description}
  \item[Network Probe DaemonSet.]
    A custom network monitoring agent
    (\texttt{mncc-underlay-network-probe:v1.0.0}) is deployed on
    all five worker nodes as a Kubernetes DaemonSet.
    Each probe instance performs active measurements of underlay
    link quality, reporting four metric classes: \emph{throughput},
    \emph{latency}, \emph{packet loss}, and \emph{link failure
    detection}.
    The measurement cycle operates with a configurable inter-probe
    delay of \textbf{120~seconds}, yielding a sampling frequency of
    $f_s = 1/120 \approx 8.3 \times 10^{-3}$~Hz per link per
    metric.
    Metrics are exposed via a Prometheus-compatible endpoint on
    port~5001.

  \item[NeMeX (Network Metrics Exporter).]
    The NeMeX component (\texttt{mncc-nemex:latest}) aggregates
    network performance metrics from the probe DaemonSet and
    computes derived indicators (e.g., link health scores,
    threshold violations) used by the mNCC self-healing logic for
    fault classification and remediation triggering.

  \item[Prometheus \& Grafana.]
    A full Prometheus Operator deployment
    (namespace: \texttt{nemo-net}) scrapes all instrumented
    endpoints at 15-second intervals, providing long-term metric
    storage with a default retention of 15~days.
    Grafana dashboards visualise real-time and historical cluster
    state.
    A secondary Prometheus instance (namespace: \texttt{nemo-ppef})
    with node-exporter DaemonSets provides host-level CPU, memory,
    disk, and network utilisation metrics.

  \item[Kepler (energy).]
    The Kepler DaemonSet (v0.7.10) provides per-pod and per-node
    energy consumption estimates derived from hardware performance
    counters (RAPL) and eBPF-based resource accounting, enabling
    correlation of provisioning overhead with energy cost.

  \item[Falco (security audit).]
    Falco is deployed on all six nodes (namespace: \texttt{nemo-svc})
    for runtime security monitoring and syscall-level audit logging,
    relevant to the security overhead evaluation in
    Section~\ref{subsec:security_methodology}.
\end{description}

\subsubsection{mNCC Control Plane Architecture}
\label{subsubsec:mncc_controlplane}

The mNCC control plane components relevant to the evaluation are
deployed in the \texttt{nemo-net} namespace:

\begin{enumerate}
  \item \textbf{Intent-Based System (IBS)} ---
    Receives network intents via the RabbitMQ message broker
    (exchange: \texttt{nemo.api.workload}, routing key:
    \texttt{intent-notify}).
    The IBS parses the intent JSON payload, resolves the
    \texttt{objectType} (e.g., \texttt{L2SM\_NETWORK}), and
    dispatches CRUD operations to the appropriate Connectivity
    Controller via gRPC.
    Responses are published to the \texttt{mncc} exchange
    (routing key: \texttt{mncc.ibs}).

  \item \textbf{L2S-M Connectivity Controller} ---
    Implements the network provisioning logic: overlay segment
    creation, NED configuration, tunnel establishment, and
    security policy enforcement (MACsec, RBAC).
    Exposes a gRPC API on port~50051.

  \item \textbf{Identity and Credential Provider} ---
    The \texttt{l2sm-idcoprovider} manages ONOS controller
    credentials and SDN southbound authentication.

  \item \textbf{Security Gateway (Kong + OAuth2)} ---
    Deployed in \texttt{nemo-sec}, the Kong API gateway with
    OAuth2 proxy enforces authentication and rate limiting on
    external API access.
    RabbitMQ handles inter-component message authentication via
    AMQP credentials.
\end{enumerate}

\subsubsection{Measurement Tools}
\label{subsubsec:tools}

Table~\ref{tab:tools} summarises the measurement tools employed
across the five experiments.

\begin{table}[htbp]
  \centering
  \caption{Measurement tools and their configuration per experiment.}
  \label{tab:tools}
  \resizebox{\columnwidth}{!}{%
  \begin{tabular}{lllc}
    \toprule
    \textbf{Experiment} & \textbf{Tool} & \textbf{Configuration} & \textbf{Sampling} \\
    \midrule
    E1--E2 (Provisioning)
      & \texttt{clock\_gettime} & \texttt{CLOCK\_MONOTONIC}, ns precision & per-event \\
      & RabbitMQ timestamps     & AMQP message headers                    & per-intent \\
    \midrule
    E3 (Self-healing)
      & UDP probe stream & 100-ms interval, custom sender/receiver & 10~Hz \\
      & NeMeX fault log  & Threshold-based event classification     & per-event \\
      & \texttt{kubectl} & Node drain, cordon, uncordon             & per-trial \\
    \midrule
    E4 (Security CP)
      & \texttt{clock\_gettime} & Per-security-layer instrumentation & per-request \\
      & \texttt{kubectl auth}   & RBAC evaluation via API server     & per-request \\
    \midrule
    E5 (Data plane)
      & \texttt{iperf3}  & 8 parallel streams, 60-s window        & per-trial \\
      & ICMP ping        & RTT/2 approximation, 100 probes        & per-trial \\
      & Metrics API      & Node CPU utilisation                   & per-trial \\
    \bottomrule
  \end{tabular}%
  }
\end{table}

\subsubsection{Baseline Characterisation}
\label{subsubsec:baseline_char}

Prior to the controlled experiments, a 24-hour monitoring campaign
was conducted to characterise the testbed's steady-state behaviour.
The network probe DaemonSet recorded underlay metrics across all
$\binom{5}{2} = 10$ worker-to-worker link pairs at 120-second
intervals, yielding 720 samples per link over the observation
window.
Table~\ref{tab:stats} reports the aggregate statistics.

\begin{table}[htbp]
  \centering
  \caption{Underlay network baseline statistics (24-hour campaign).}
  \label{tab:stats}
  \begin{tabular}{lcccc}
    \toprule
    \textbf{Metric} & \textbf{Mean} & \textbf{Std.} & \textbf{CV} & \textbf{P99} \\
    \midrule
    Throughput (Gbps)     & $\sim$9.4  & 1.3  & 0.14 & 7.1 \\
    One-way latency (ms)  & $\sim$0.18 & 0.03 & 0.17 & 0.28 \\
    Packet loss (\%)      & $<$0.01    & ---  & ---  & 0.02 \\
    \bottomrule
  \end{tabular}
\end{table}

The coefficient of variation $\mathrm{CV} \approx 0.14$ for
throughput confirms low temporal variability, validating the
power analysis assumptions used to determine the sample sizes
in Section~\ref{subsec:methodology}.
The sub-millisecond baseline latency and negligible packet loss
indicate a healthy underlay suitable for isolating overlay
provisioning effects from infrastructure noise.

\subsubsection{Reproducibility}
\label{subsubsec:reproducibility}

The complete testbed configuration---including Kubernetes manifests,
Helm charts, network probe deployment, mNCC component images, and
all measurement scripts---is archived in the project repository.
The cluster state can be restored to the evaluation baseline by
applying the versioned manifests against a fresh Kubernetes~v1.29
installation on equivalent hardware.
All raw measurement data, statistical analysis notebooks, and
plotting scripts are publicly available to enable independent
verification of the results presented in
Section~\ref{sec:results}.
